\section{Lifting the obstruction to a line bundle}\label{sec:bundle}

Proposition~\ref{prop:coordinate-obstruction} reduces noncoordinateness to
excluding a nonzero locally nilpotent derivation $E$ of $B$ with a nonzero
invariant in $I=(f_0,g_0)$.  We pull $I$ back to the complement of the zero
section in a line bundle over $\Spec B$; on this complement the ideal
becomes principal.  Every additive action on the base lifts linearly to
the line bundle and preserves this complement.  Consequently, an invariant
in $I$ will force the lifted action to fix one specific polynomial.  This use of the quadric and its line bundle
adapts the construction of \cite[Sections~4--6]{Cancellation2026}; all the
lifting and principalization arguments needed here are proved below.

Put $Y=\Spec B$ and define
\begin{equation}\label{bun:C}
 C=k[a,d,b,c,u]/(ac-bd-1).
\end{equation}
The assignment
\begin{equation}\label{bun:base-map}
 x=ab,\qquad y=dc,\qquad z=db,\qquad u=u
\end{equation}
defines a homomorphism $B\longrightarrow C$, since
$xy=abdc=(db)(ac)=z(z+1)$.  The ring $C$ is a domain: the polynomial
$ac-bd-1$, viewed as a primitive linear polynomial in $c$ over
$k[a,d,b,u]$, is irreducible.  We give $C$ the \emph{fiber weights}
\begin{equation}\label{bun:fiber-weights}
 \operatorname{wt}(a,d,b,c,u)=(1,1,-1,-1,0).
\end{equation}
The relation is homogeneous of weight zero, and the image of $B$ has weight
zero.  These weights describe the action of $\G_m$ that scales the fibers
of the following bundle.  Geometrically this is the familiar quotient of
$\mathrm{SL}_2$ by a diagonal torus, times the affine $u$-line; see
\cite[p.~8431]{BlancVanSanten2019}.  We use explicit charts throughout.

\begin{lemma}\label{bun:charts}
The map $B\longrightarrow C$ is injective.  The morphism
$\Spec C\longrightarrow Y$ is the complement of the zero section in a line
bundle $\mathcal L\longrightarrow Y$.  It is trivial over the open cover
\[
 U_0=D(z+1),\qquad U_1=D(z)
\]
of $Y$, with respective invertible fiber coordinates $a$ and $d$.
\end{lemma}

\begin{proof}
Over $U_0$, the identity $ac=z+1$ makes $a$ invertible.  The formulas
\begin{equation}\label{bun:chart-a}
 c=\frac{z+1}{a},\qquad b=\frac{x}{a},\qquad
 d=\frac{ay}{z+1}
\end{equation}
give an isomorphism
$C[(z+1)^{-1}]=B[(z+1)^{-1}][a,a^{-1}]$.
In particular the localized map from $B$ is injective, and so is the
original map because $B$ is a domain.  Over $U_1$, the corresponding formulas
are
\begin{equation}\label{bun:chart-d}
 b=\frac{z}{d},\qquad c=\frac{y}{d},\qquad a=\frac{xd}{z},
\end{equation}
and give $C[z^{-1}]=B[z^{-1}][d,d^{-1}]$.
The opens cover $Y$ because $z$ and $z+1$ generate the unit ideal.  On their
overlap the fiber coordinates satisfy
\begin{equation}\label{bun:transition}
 d=\frac{y}{z+1}\,a,
\end{equation}
where $y/(z+1)$ is a unit with inverse $x/z$.  Gluing the trivial line bundles
by this transition gives $\mathcal L$; removing their zero sections gives
precisely the two charts of $\Spec C$ above.  Both fiber coordinates have
weight one, so the scaling agrees with \eqref{bun:fiber-weights}.
\end{proof}

The two generators of $I$ acquire a common factor in $C$.  Define
\begin{equation}\label{bun:v}
 v=a^2(x+u^2)-d^2=a^3b+a^2u^2-d^2\in C.
\end{equation}

\begin{lemma}\label{bun:principalization}
The extended ideal $IC$ is the principal ideal $(v)$.
\end{lemma}

\begin{proof}
In $C[a^{-1}]$, put $\xi=d/a$ and $j=x+u^2-\xi^2$.  The determinant relation
and \eqref{bun:base-map} give
\[
 z=\xi x,\qquad y=\xi+\xi^2x,\qquad
 s_B=y-x(x+u^2)=\xi-xj.
\]
Consequently
\begin{equation}\label{bun:ideal-chart}
 f_0=(1+2\xi x-x^2j)j,\qquad g_0=x^2j,\qquad v=a^2j.
\end{equation}
The two coefficients of $j$ generate the unit ideal, as the identity
\[
 (1-2\xi x)(1+2\xi x-x^2j)
   +x^2(4\xi^2+j-2\xi xj)=1
\]
shows.  Thus $IC[a^{-1}]=(j)=(v)$.

It remains to check primes of $C$ containing $a$.  Modulo such a prime,
$bd=-1$, while $x=0$ and $z=-1$.  Hence $g_0=z-s_Bx$ has nonzero residue
$-1$, and $v$ has nonzero residue $-d^2$.  Both are units in the corresponding
local ring.  The ideals $IC$ and $(v)$ therefore agree at every prime, which
proves the assertion.
\end{proof}

Figure~\ref{bun:figure} records the geometric meaning of this ideal
equality: the inverse image of $V(I)$ is cut out by the single function $v$.

\begin{figure}[ht]
\centering
\begin{tikzcd}[column sep=large,row sep=large]
 V(v) \arrow[r,hook] \arrow[d] & \Spec C \arrow[d] \\
 V(I) \arrow[r,hook] & Y=\Spec B
\end{tikzcd}
\caption{A Cartesian square of closed subschemes.  Pulling $I$ to the
complement of the zero section of $\mathcal L$ gives the principal ideal
$(v)$.  In particular, every $h\in I$ pulls back to a multiple of $v$.}
\label{bun:figure}
\end{figure}

We next lift additive actions from $Y$ to $\mathcal L$.  The elementary
reason that a lift exists is that adjoining an affine-line parameter does
not change line bundles on $Y$.

\begin{lemma}\label{bun:homotopy}
Every line bundle on $\A^1\times Y$ is isomorphic to the pullback of its
restriction to $\{0\}\times Y$.
\end{lemma}

\begin{proof}
The two charts of Lemma~\ref{bun:charts} have coordinate rings
\begin{align*}
 B[(z+1)^{-1}]&=k[x,t,u,(xt+1)^{-1}],
       &t&=\frac{y}{z+1},\quad z=xt,\\
 B[z^{-1}]&=k[x,t,u,(xt-1)^{-1}],
       &t&=\frac{y}{z},\quad z=xt-1.
\end{align*}
Here $t$ denotes a separate coordinate on each chart.  Each ring, and its
polynomial extension by the affine-line coordinate $T$, is a noetherian
unique factorization domain.  A line bundle on the spectrum of such a
ring is trivial: its divisor class is represented by a Cartier divisor,
and unique factorization makes every divisor principal.

Let $\mathcal M$ be a line bundle on $\A^1\times Y$.  Choose trivializations
on $\A^1\times U_0$ and $\A^1\times U_1$.  Its transition function belongs
to $B[z^{-1},(z+1)^{-1}][T]^*$.  A polynomial over a domain is a unit only
when it is a constant unit, so this transition function is independent of
$T$.  Its specialization at $T=0$ is therefore the same transition
function.  The chosen trivializations identify $\mathcal M$ with the
pullback of its restriction at zero.
\end{proof}

The lifting argument has a useful general form.  Its normalization step
is the standard one for linearizing a line bundle under an additive action;
compare \cite[Lemma~2.9]{Brion2015}.

\begin{lemma}\label{bun:action-lift}
Let $X$ be an affine integral $k$-variety such that every line bundle on
$\A^1\times X$ is the pullback of its restriction at zero.  Every
$\G_a$-action on $X$ lifts to any line bundle on $X$ by maps linear on
fibers, commuting with fiber scaling.
\end{lemma}

\begin{proof}
Write $R=k[X]$, let $\mathcal M$ be the line bundle, and let
$\alpha:\A^1\times X\longrightarrow X$ be the action.  If
$\pi:\A^1\times X\longrightarrow X$ is projection, then
$\alpha^*\mathcal M$ and $\pi^*\mathcal M$ restrict to $\mathcal M$ at
parameter zero.  The hypothesis therefore supplies an isomorphism
\[
 \eta:\pi^*\mathcal M\longrightarrow\alpha^*\mathcal M.
\]
At zero, $\eta$ is an automorphism of $\mathcal M$, hence multiplication
by a unit of $R$.  Multiplying $\eta$ by the inverse pullback of that unit
normalizes it to the identity at zero.

The isomorphism $\eta$ gives a fiberwise linear lift of $\alpha$.  To check
the group law, compare, over $\A^2\times X$, the map on fibers obtained
by applying parameters $S$ and then $T$ with the map obtained by applying
$T+S$.  Both are global isomorphisms of the same source and target line bundles,
with fiber maps from $\mathcal M_x$ to
$\mathcal M_{\alpha(T+S,x)}$.  Their ratio is an automorphism of the source
line bundle.  Its endomorphism sheaf is canonically $\mathcal O$, so this
automorphism is multiplication by a global unit in $R[T,S]$.
Since $R$ is a domain, this unit is independent of $T,S$.  At $(T,S)=(0,0)$
it equals one by the normalization of $\eta$.  Thus the group law holds.
The resulting action is linear on fibers, so it commutes with fiber
scaling and preserves the complement of the zero section.
\end{proof}

We will also use the following elementary property of locally nilpotent
derivations; see \cite[Proposition~2.1(3)]{KalimanMakarLimanov2007}.

\begin{lemma}\label{bun:factorial-closure}
Let $R$ be a domain over a field of characteristic zero and let $\delta$ be a locally
nilpotent derivation of $R$.  If $r,s\in R$ are nonzero and $\delta(rs)=0$,
then $\delta(r)=\delta(s)=0$.
\end{lemma}

\begin{proof}
The exponential $\exp(T\delta)$ takes values in $R[T]$.  If $rs$ is fixed,
then
\[
 \exp(T\delta)(r)\,\exp(T\delta)(s)=rs.
\]
Degrees in $T$ add under multiplication of nonzero polynomials over a
domain.  Both factors on the left therefore have degree zero, and their
coefficients of $T$ give the assertion.
\end{proof}

\begin{proposition}\label{prop:lift}
Let $E$ be a nonzero locally nilpotent derivation of $B$ such that
$I\cap\ker E\ne0$.  There is a nonzero locally nilpotent derivation
$\widetilde E$ of $C$ that extends $E$, preserves fiber weights, and satisfies
$\widetilde E(v)=0$.
\end{proposition}

\begin{proof}
The exponential of $E$ defines an additive-group action on $Y$.
By Lemmas~\ref{bun:homotopy} and~\ref{bun:action-lift}, it lifts to
$\mathcal L$ by fiberwise linear maps.  Restrict to the complement of the
zero section, namely $\Spec C$.
Differentiating this action at the identity gives a derivation
$\widetilde E$ of $C$.  The action law expresses its pullback on each
function as $\sum_{r\ge0}T^r\widetilde E^{\,r}/r!$; because the action is
algebraic, the sum is finite for each function.  Hence $\widetilde E$ is
locally nilpotent.  It extends $E$ and is nonzero, and commutation with
fiber scaling says exactly that it preserves the weights
\eqref{bun:fiber-weights}.

Finally choose $0\ne h\in I\cap\ker E$.  The inclusion $B\subset C$ and
Lemma~\ref{bun:principalization} give a factorization $h=vq$ in $C$, with
both factors nonzero.  Since $\widetilde E(h)=0$,
Lemma~\ref{bun:factorial-closure} yields $\widetilde E(v)=0$.
\end{proof}
